% Transcription — fonds Alexandre Grothendieck, shelfmark 121, batch 4 (pages 61-72).
% Established from the facsimile archives/batches/121/batch-04.pdf (a copy of
% 121.pdf fetched through the project's relay,
% https://grothendieck-relay.onrender.com/source/121.pdf; PDF page k is page
% 60+k, no cover sheet), per the skill transcribe-grothendieck. The folder is
% seventy-two pages in four batches (1-20, 21-40, 41-60, 61-72), transcribed
% in parallel; this is the last, twelve pages.
%
% Pass: Opus 5.5 (claude-opus-5-5), 2026-09-26 — first pass, unchecked against the pages by a human.
%
% Numbering: \page{} follows the PDF count, which is the Montpellier footer
% (61-72). The archivists' pencil numbers bottom left run two ahead
% throughout (63 on p. 61 ... 74 on p. 72), as batches 1 and 2 found. He
% numbers pp. 69 and 72 « 1) » and « 2) » top left, and the sections on
% them « 1. » (p. 69) and « 2. » (p. 71); nothing else is paginated.
%
% Dating: \dating{1974} only, the folder's own date; nothing on these pages
% carries a date.
%
% The letter. No letter to him appears in this batch. Page 66 carries, upside
% down on its left half, the opening of a letter draft of his own in English
% (« Dear Fujii Guruji, Please accept my most cordial wishes of welcome in
% France and of a very happy stay in Paris, »), cancelled by five diagonal
% strokes and numbered 89 in its corner; it has no mathematics and is not
% transcribed (the mathematics written around it is). Whether it is the
% « lettre (1974) » of the inventory the pages do not say. Pages 69-72
% (« 1) », « 2) ») are a fair statement of axioms, possibly written to be
% sent, but carry no address, salutation or signature; they are transcribed
% as notes.
%
% Pages skipped: 63 and 68, blank cream separator sheets.
% Pages summarised: none whole. On p. 65 the left half (turned sideways) is
% transcribed as formulas; on p. 66 the scratch computations of the left
% half are given in a \note with what can be read, the right half is
% transcribed.
% Pages transcribed: 61, 62, 64-67, 69-72; blue ink, pp. 69-72 in black.
% Pp. 64-67 are landscape double sheets: p. 64 on the blank verso of an
% Institut de Mathématiques (USTL, Montpellier) letterhead, whose printed
% half carries nothing of his; p. 67 has its right half blank.
%
% What the batch is about. P. 61 ends a « Dém » begun on p. 60 (batch 3):
% embedding X, Y, Z in cubes I_0^n and building a commutative cube of monos
% with cartesian faces, I^nu a retract of I^N, I^N', whence a map to S. P.
% 62: reformulating the data by F_n = Hom_M(I_0^n, I_0) (tame functions)
% instead of the classes M_n, with conditions (F1)-(F4) a)-c) and the graph
% criterion Gamma_f in M_{n+1}. P. 64: an axiom list for M_n in I^n, I =
% [-1,+1] (stability, simplicial maps, M-injectivity of I, triangulations,
% a linear example in M_{n+2}). P. 65: conditions a)-d) on a square X c Y,
% Z c S for S = Y u_X Z, the effective epimorphism Z u Y -> S and its
% kernel pair; sideways, Chern class formulas for Grassmannians and flag
% bundles (unrelated). Pp. 66-67: realising such a square by cubes, Y x_S Y
% = Delta_Y v (X x X). Pp. 69-72: « 1. » a set R with classes C_n of subsets
% of R^n, axioms (1)-(4), the finite-index form C_I, C-morphisms by graphs,
% C-modèles (structures by sets T_n of injections E -> R^n), the category of
% C-modèles, conditions a)-c) for a final object, (5) R^n = union of C_n,
% C-parties; « 2. » R Hausdorff, (7) members of C_n compact, (8) a
% fundamental system of neighbourhoods, the sheaf of germs of C-morphisms,
% C-espaces; « 2) » (6) stability of C_n under finite unions, gluing of
% C-morphisms on finite covers by C-parties, and the Question of amalgamated
% sums, with a condition of extension of C-morphisms Y -> R to X -> R.
%
% Notation fixed for the batch: I_0, I_0^n and I^n as he writes them (with
% and without the index 0, not regularised); \Phi, \varphi, \Psi, \psi,
% \overline{\varphi} for his barred phi; M for his script M (the category),
% \mathcal{M}_n for the classes, M_* as he writes it; F_n, F_{m,n} for his
% script F; C_n, C_I plain; his subscript mark on « C-modèle », « C-morphisme
% » (a small star or x) set as C_{*} throughout, with one note; \sqcup for
% his coproduct, \sqcup_X for the amalgamated sum; e for the final object;
% circled axiom numbers written (n) with a note.
%
% Readings to check first: the top lines of p. 61 (« qui ... deux sections
% »), the struck passage at the foot of p. 62, the struck and boxed lines of
% p. 64, the letters d), a), b), c) of p. 65, the diagonal margins of pp. 69
% and 70, the c) of p. 69, the last lines of the Prop. on p. 70, the margin
% of p. 72.
% Apparatus: 89 \ill{} and 20 \uncertain{} (counted from the source,
% including those inside \struck{}, \note{} and \marginal{}).

\documentclass[11pt,a4paper]{article}
% Le préambule commun à toutes les transcriptions du fonds Grothendieck.
%
% Il tient en une idée : **l'incertitude se compose, elle ne se résout pas.**
% Une transcription qui lisse un mot douteux a détruit la seule chose pour
% laquelle on la faisait — un lecteur doit pouvoir savoir, à chaque endroit,
% ce qui était sur le feuillet et ce qui a été ajouté. Les six macros qui
% suivent sont donc l'appareil critique, et non de la décoration.
%
% Les mêmes commandes sont reconnues par `scripts/render.mjs`, qui produit la
% vue de lecture du site. Une macro ajoutée ici doit l'être aussi là-bas,
% faute de quoi le rendu échouera bruyamment — ce qui est voulu : un rendu
% silencieusement faux est pire qu'un rendu absent.

\ProvidesPackage{grothendieck}[2026/08/08 Transcriptions du fonds Grothendieck]

\usepackage{amsmath,amssymb,amsthm}
\usepackage{mathrsfs}
\usepackage{stmaryrd}
% Les diagrammes commutatifs sont partout dans ce fonds : c'est la langue
% même de la géométrie algébrique de Grothendieck, et une transcription qui
% les rendrait en prose ne transcrirait rien.
\usepackage{tikz-cd}
% Français par défaut — c'est la langue des feuillets — mais l'anglais est
% chargé aussi : la traduction et le résumé sont des documents anglais, et la
% typographie française (espaces avant les deux-points) y serait une faute.
% Ces documents basculent par \selectlanguage{english} en tête de corps.
\usepackage[english,french]{babel}
\usepackage{fontspec}
\usepackage{geometry}
\geometry{a4paper,margin=25mm,marginparwidth=28mm}
\usepackage{marginnote}
\usepackage{xcolor}
% ulem, not soul. `soul`'s \st and \ul are built on a character-by-character
% scan that XeLaTeX's font machinery breaks: the document fails with an
% undefined control sequence rather than degrading. ulem does the same two jobs
% and is Unicode-engine safe. `normalem` stops it hijacking \emph, which the
% transcriptions use for Grothendieck's own emphasis.
\usepackage[normalem]{ulem}
\usepackage{hyperref}
\hypersetup{colorlinks=true,linkcolor=black,urlcolor=black,pdfborder={0 0 0}}

% --- Métadonnées du lot -----------------------------------------------------
% Reprises telles quelles par le script de rendu, qui les affiche en tête de
% la vue de lecture. Elles fixent ce que le fichier prétend couvrir : sans
% elles, une transcription flotte sans point d'attache dans un fonds de seize
% mille feuillets.

\newcommand{\folder}[1]{\gdef\grothFolder{#1}}
\newcommand{\batch}[1]{\gdef\grothBatch{#1}}
\newcommand{\leaves}[2]{\gdef\grothFirst{#1}\gdef\grothLast{#2}}
\newcommand{\foldertitle}[1]{\gdef\grothFoldertitle{#1}}
\gdef\grothFolder{?}\gdef\grothBatch{?}\gdef\grothFirst{?}\gdef\grothLast{?}\gdef\grothFoldertitle{}

% --- Le résumé --------------------------------------------------------------
%
% Il ouvre la lecture modernisée, sous le titre « Résumé », et il est composé
% en petit. Ce n'est pas un ornement : c'est le seul passage du document qui
% ne soit pas des mathématiques, et un lecteur doit pouvoir le reconnaître
% avant de l'avoir lu — puis le sauter s'il connaît déjà le sujet, ou s'y
% arrêter s'il ne le connaît pas. Le filet qui le suit marque la reprise du
% corps.

\newenvironment{resume}
  {\small\setlength{\parskip}{0.45em}}
  {\par\medskip\hrule\medskip}

% --- Mots-clés ---------------------------------------------------------------
%
% En anglais, en clôture du résumé de la lecture modernisée : le vocabulaire
% moderne sous lequel chercher ce que les feuillets construisent. Le manifeste
% du site (`npm run manifest`) les extrait pour étiqueter la cote entière —
% cette ligne est la source unique des tags, il n'y a pas de fichier à côté.
\newcommand{\keywords}[1]{\par\smallskip\noindent
  {\footnotesize\textsc{Keywords} --- \textit{#1}}\par}

% --- L'appareil critique ----------------------------------------------------

% Le numéro de feuillet, en marge. C'est l'ancre qui permet de régler un
% désaccord sur une formule en regardant *un* feuillet plutôt qu'une chemise
% de six cents. Le site s'en sert aussi pour tourner les pages du fac-similé
% quand on fait défiler la transcription.
\newcommand{\leaf}[1]{%
  \marginnote{\footnotesize\color{gray!70}\textsf{#1}}%
  \label{leaf:#1}\ignorespaces}

% Illisible : on ne devine pas. Un mot inventé qui se lit comme les autres est
% le pire résultat possible.
\newcommand{\ill}{\textcolor{red!60!black}{[\,\ldots\,]}}

% Lecture douteuse : le mot est donné, et signalé comme incertain.
\newcommand{\uncertain}[1]{\uline{#1}}

% Ajout de l'éditeur — une lettre manquante, un mot sous-entendu.
\newcommand{\add}[1]{\textcolor{blue!50!black}{[#1]}}

% Biffé par Grothendieck lui-même. Ce qu'il a rayé fait partie du document :
% une rature dit souvent où la pensée a changé de direction.
\newcommand{\struck}[1]{\sout{#1}}

% Note du transcripteur, sur l'état matériel du feuillet.
\newcommand{\note}[1]{\par\smallskip{\small\color{gray!60!black}\textit{[#1]}}\par\smallskip}

% Note marginale de Grothendieck, distincte du corps du texte.
\newcommand{\marginal}[1]{\par\smallskip\noindent\hspace{1em}%
  {\small\color{gray!60!black}#1}\par\smallskip}

% --- Titre ------------------------------------------------------------------

\AtBeginDocument{%
  \begin{center}
    {\large\bfseries Fonds Alexandre Grothendieck --- cote n\textsuperscript{o}~\grothFolder}\\[2pt]
    {\itshape\grothFoldertitle}\\[4pt]
    {\small feuillets \grothFirst--\grothLast\ (lot \grothBatch)}\\[2pt]
    {\footnotesize\color{gray!60!black}Transcription établie d'après le fac-similé numérisé
     par l'Université de Montpellier.\\
     Les lectures incertaines sont soulignées, les illisibles notées [\,\ldots\,],
     les ajouts entre crochets.}
  \end{center}
  \vspace{1em}}

\endinput


\folder{121}
\batch{4}
\pages{61}{72}
\dating{1974}
\watermark{Édition de démonstration}
\foldertitle{Topologie modérée : notes manuscrites (s.d.), lettre (1974)}

\begin{document}

\page{61}
\emph{Dém}\note{suite de la page 60, où sont posées les conditions 5°)
et 6°)} Si $X$ est rétract dans $Y$, $Z$ \add{(dans $M$)}, alors on plonge
\struck{$Y$ dans} $Y$ et $Z$ dans $Y \times_X Z$\note{sic : le signe est
un produit fibré ; on attendrait plutôt une somme amalgamée} qui \ill{}
\uncertain{deux sections} $X \to Y$, $X \to Z$, et \ill{} \ill{} \ill{}
$Y \cap Z = X$, OK. \struck{Il faudrait} On se ramène au cas particulier.
Partant d'un diagramme cartésien (en 6°)

\begin{tikzcd}
X \arrow[r, hook] \arrow[d] & Y \arrow[d, "\Psi"] \\
e \arrow[r] & I_0^p
\end{tikzcd}

et considérons une $Y \overset{\Phi}{\hookrightarrow} I_0^n$ mono, $\Phi
| X = \varphi$, d'où \struck{$Y \to$} diagramme cartésien de monos

\begin{tikzcd}
X \arrow[r, hook] \arrow[d, "{(e,\ \varphi)}"'] & Y \arrow[d, "{(\Psi \times \Phi)}"] \\
e \times I_0^n \simeq I_0^n \arrow[r, hook] & I_0^p \times I_0^n \simeq I_0^{p+n}
\end{tikzcd}

\note{l'étiquette de gauche est surchargée : le premier terme, biffé, se
lit mal}

et de même

\begin{tikzcd}
X \arrow[r, hook] \arrow[d, "\varphi'"'] & Z \arrow[d, "{(\Psi',\ \Phi')}"] \\
e \times I_0^{n'} \simeq I_0^{n'} \arrow[r, hook] & I_0^{p'} \times I_0^{n'} = I_0^{p'+n'}
\end{tikzcd}

Mais quitte à prolonger $\varphi$ de $X$ à $Z$ en $\overline{\varphi} : Z
\to I_0^n$, et $\varphi'$ \struck{de $X$} à $Y$ en $\overline{\varphi}' : Y
\to I_0^{n'}$, on trouve des diagrammes cart.

\begin{tikzcd}
X \arrow[r, hook] \arrow[d, "{(\varphi',\ \varphi)}"'] & Z \arrow[d, "{(\Psi',\ \Phi',\ \overline{\varphi})}"] \\
e \times I_0^{n'} \times I_0^n \arrow[r, hook] & I_0^{p'} \times I_0^{n'} \times I_0^n
\end{tikzcd}

\begin{tikzcd}
X \arrow[r, hook] \arrow[d, "{(\varphi,\ \varphi')}"'] & Y \arrow[d, "{(\Psi,\ \Phi,\ \overline{\varphi}')}"] \\
e \times I_0^n \times I_0^{n'} \arrow[r, hook] & I_0^p \times I_0^n \times I_0^{n'}
\end{tikzcd}

\[
  n + n' = \nu, \qquad p + n + n' = N, \qquad p' + n + n' = N'
\]
\note{ces trois égalités sont écrites en marge droite, à hauteur du second
diagramme}

et finalement un diagramme \add{comm.} de monos

\begin{tikzcd}[column sep=small, row sep=small]
 & Z \arrow[dd] & & \\
X \arrow[ur, hook] \arrow[rr, hook] \arrow[dd] & & Y \arrow[dd] & \\
 & I^{N'} \arrow[rr, dashed] & & S \\
I^{\nu} \arrow[ur, hook] \arrow[rr, hook] & & I^{N} \arrow[ur, dashed] &
\end{tikzcd}

\note{cube redessiné d'après le croquis : les deux flèches vers $S$ sont en
pointillé ; l'emplacement exact des flèches verticales issues de $Z$ et de
$Y$ est serré sur la page}

avec carrés cart., et $I^{\nu}$ rétract de $I^N$, $I^{N'}$. Dès lors il y a
$\dashrightarrow S$ donnant carré cartésien. OK

\page{62}
Il ne semble pas que 5° et 6° s'expriment de façon simple en termes des
$\mathcal{M}_n$. Mais \ill{} \ill{} \ill{} \uncertain{profiter} pour
reformuler les données ainsi. Pour tt $n \geqslant 0$, soit
\[
  F_n = \mathrm{Hom}_M(I^n, I) \wr \ \mathrm{Hom}_M(I_0^n, I_0)
\]
\note{il écrit « $F_n = \mathrm{Hom}(I^n, I)$ », un signe $\wr$ vertical
dessous, puis « l'ens.\ $\mathrm{Hom}_M(I_0^n, I_0)$ » ; un mot biffé à
droite de la première formule}
l'ens.\ $\mathrm{Hom}_M(I_0^n, I_0)$. Grâce à 6°), la connaissance des
$F_n$ implique celle des $\mathcal{M}_n$. Quelles conditions \ill{} ?

\struck{(F1)} Posons $F_{m,n} \subset \mathrm{Hom}_K(I^n, I^m)$ défini
(composantes par composantes) par $\varphi = (\varphi_1, \ldots,
\varphi_m)$ $F$-modéré ssi les $\varphi_i \in F_n$, \struck{\ill{}} \ill{}
\begin{itemize}
  \item[(F1)] $F_n$ contient les $\mathrm{pr}_i : I^n \to I$.
  \item[(F2)] Stabilité par composition des $F_{m,n}$.
\end{itemize}
\note{« (F1) » est d'abord écrit devant « Posons », puis biffé et reporté,
serré, à la ligne suivante ; devant (F2), un signe biffé}

On pose $\mathcal{M}_n =$ ens.\ des \ill{} objets de $I^n$
\add{intersections finies de sous-objets \struck{\ill{}}} de la forme
$\varphi^{-1}(\alpha)$, où $\varphi \in F_n$, $\alpha \in F_0 : e
\overset{\alpha}{\hookrightarrow} I$.

Alors $\mathcal{M}_{*}$ satisfait M1, M2, pour avoir M3 \uncertain{bis}
(\ill{} \ill{}) on doit admettre
\begin{itemize}
  \item[(F3)] Si $\varphi_1, \ldots, \varphi_m \in F_n$, d'où $I^n \to I^m$
  \ldots\note{sous « (F3) », « (M3 bis) » et, à la suite, « et si » ; la
  ligne s'achève sur des points de suspension}
  \item[(F4)] a) $\exists\, 0 : e \to I$ \add{$\in F_0$} tel que $\forall
  \alpha \in F_1$, $\alpha : e \to I$, on ait $\mathrm{Im}\, \alpha =
  \bigcap_{1 \leqslant i \leqslant n} \varphi_i^{-1}(0)$, pour $\varphi_1,
  \ldots, \varphi_n \in F_1$ \quad ($\mathcal{M}_1^{\cap}$ i.e.)\note{le
  chiffre de (F4) est surchargé ; « $\forall \alpha \in F_1$ » est écrit
  tel quel, bien que $\alpha$ soit pris dans $F_0$ à la ligne précédente}

  b) \struck{\ill{}} $(e \times I \,;\, I \times e) \subset I^2$ est $\in
  \mathcal{M}_2$, i.e.\ est de la forme $\bigcap_i \varphi_i^{-1}(0)$, où
  $\varphi_i \in F_2$.

  c) $\exists\, \varepsilon_1, \varepsilon_2 \in F_0$, $\varepsilon_1,
  \varepsilon_2 : e \to I$, tel que $e \sqcup e \to I$ soit mono.
\end{itemize}

Il \struck{\ill{} exprimer \ill{} que pour $f : X \to I$ \ill{} diagramme
$\Gamma_f$}\note{passage biffé et récrit au-dessus de lui-même, à peine
lisible ; on y lit encore « $\subset I^n \times$ »} $\Gamma_f \subset
I^{n+1}$, \struck{$F_f$} \ill{} pour que $f \in F_n$, il faut et il suffit
que $\Gamma_f \in \mathcal{M}_{n+1}$ ; \struck{a) Si $f \in F_n$,
$\exists$} le « il faut » \ill{} i.e.\ $\Delta_I \subset I \times I$ est $\in
\mathcal{M}_2$ ; le « il suffit » doit être postulé indépendamment.

\page{64}
\[
  \mathcal{M}_n \subset I^n \qquad n \in \mathbb{N} \qquad (I = [-1, +1])
\]

\begin{itemize}
  \item[] $\mathcal{M}_n$ stable par int.\ et réunions finies
  \item[] $\mathcal{M}_{*}$ stables par images directes et inverses
  d'appl.\ \uncertain{simpliciales}
  \item[] $\mathcal{M}_0$ contient \struck{deux parties \ill{} \ill{}}
  $\lbrace 0 \rbrace$ et $\lbrace 1 \rbrace$
  \item[] \struck{\ill{}} [ $I$ est $M$-injectif (\ill{} \ill{} \ill{}
  \ill{} de $\mathcal{M}_{*}$) ; Tout \struck{\ill{}} $\mathcal{M}_n$ est
  de la forme $f^{-1}(0) = \bigcap_{1 \leqslant i \leqslant p}
  f_i^{-1}(0)$, $f : I^n \to I^p$ « permis ») ]\note{les deux lignes sont
  réunies par une accolade dans un crochet ; une double flèche, dans la
  marge, pointe vers elles}
\end{itemize}

a) $\mathcal{M}_{*}$ contient les simplexes types et \struck{\ill{} \ill{}}
\ill{} \uncertain{linéaires} d'un simplexe type dans $I$ \ill{} un
$M$-\struck{\ill{}} \ill{} ses sommets aux sommets \ill{}\note{ligne
surchargée de biffures et d'ajouts interlinéaires ; la lettre a) est
écrite dans la marge}

b) Pour tout « schéma simplicial » fini $\Sigma$, sa réalisation
\struck{Soit $M^a$ la catégorie modérée des \ill{}} \add{[objets \ill{}
par morceaux.]}\note{addition encadrée sous la ligne biffée} géométrique
$\widetilde{\Sigma}$ admet une $M$-structure admissible (définie par une
\emph{liste} finie dans $M$), et tout objet de $M$ [i.e.\ \struck{\ill{}}
\ill{} d'un $\mathcal{M}_n$] admet une telle « triangulation ». Pour tout
couple $Y \subset X$ dans $M$ (OPS $Y, X \in \mathcal{M}_n$, $Y \subset X$),
\ill{} \ill{} $M$-triangulation de $Y$ est induite par une
$M$-triangulation de $X$.\note{le paragraphe b) est encadré d'un trait
vertical dans la marge}

\[
  \Delta \quad \bigl\lbrace (x_0, \ldots, x_n, y) \bigm| x_i \geqslant 0,\ \
  \textstyle\sum x_i = 1,\ \ y - \sum_{0 \leqslant i \leqslant n} a_i x_i
  = 0 \bigr\rbrace \in \mathcal{M}_{n+2}
\]
si les $\lbrace a_i \rbrace \in \mathcal{M}_1$\note{$\Delta$ est un petit
triangle dessiné ; les deux conditions sont superposées dans une accolade ;
le signe entre $y$ et $\Sigma$ est un gros point, lu comme un signe moins}

\page{65}
\note{feuillet double à l'italienne ; la moitié gauche est écrite en
travers, la moitié droite droit ; on donne la droite d'abord}

\begin{tikzcd}
X \arrow[r, hook, "i"] \arrow[d, "p"'] & Y \arrow[d, "p'"] \\
Z \arrow[r, hook, "i'"'] & S
\end{tikzcd}

comm.

\begin{itemize}
  \item[d)] $S = \mathrm{Im}\, i' \cup \mathrm{Im}\, p'$ \quad i.e.\ \quad
  $Z \sqcup Y \to S$ épi effectif\note{« effectif » est ajouté serré à
  droite ; un petit signe en exposant de $S$}
  \item[a)] $i'$ mono
  \item[b)] carré cartésien i.e.\ $X = p'^{-1}(Z)$
  \item[c)] $Y \times_S Y = \Delta_Y \cup (X \times_Z X)$
\end{itemize}
\note{les lettres d) et a) sont repassées sur d'autres ; celle de c) est
écrite sur un d)}
\[
  \Longrightarrow \quad S \xleftarrow{\ \sim\ } Y \sqcup_X Z
\]
\[
  (Z \sqcup Y) \times_S (Z \sqcup Y) = Z \sqcup X \sqcup X \sqcup (Y
  \times_S Y)
\]
\[
  Z \sqcup X \sqcup X \sqcup (Y \times_S Y) \rightrightarrows Z \sqcup Y
  \to S
\]

\[
  Z \sqcup X \sqcup X \sqcup \bigl[ (X \times_Z X) \sqcup \Delta_Y \bigr]
  \rightrightarrows Z \sqcup Y
\]

\begin{tikzcd}
Z \sqcup Y \arrow[r] \arrow[dr, "{(f,\ g)}"'] & S \arrow[d, dashed] \\
 & T
\end{tikzcd}

\note{sous la dernière ligne, « $Z \sqcup Y \to S$ » récrit et biffé ; la double flèche est suivie, sur la page, de $Z \sqcup Y \to S$ et de la flèche $(f, g)$ vers $T$, redessinées à part ; en
bas de la page, deux croquis : un rectangle partagé en deux par un trait
vertical, et deux triangles opposés par un sommet}

\note{moitié gauche, écrite en travers :}
\[
  \begin{array}{c} E \\ \mid \\ S \end{array} \qquad
  \mathrm{Grass}_n(E) \leftarrow \Sigma_{n,\sigma}(E) \to D_\sigma(E),
  \qquad \mathrm{Drap}(N) \to \mathrm{Grass}_n(E)
\]
\note{diagramme redessiné en ligne : $\Sigma_{n,\sigma}(E)$ est au centre,
$\mathrm{Drap}(N)$ dessous, avec une flèche vers $D_\sigma(E)$ surchargée ;
des flèches longues vont de $\mathrm{Grass}_n(E)$ et de $D_\sigma(E)$ vers
$S$}
\[
  \begin{array}{ccc} E & & E' \\ \mid & & \mid \\ S & \text{---} & S' =
  D_\sigma(E) \end{array}
\]
\[
  H(S)\bigl[ c_i(M), c_j(N) \bigr] \Big/ \Bigl( \textstyle\sum
  c_i(M) \Bigr) \Bigl( \sum c_j(N) \Bigr) = \sum_{h} c_h(E)
\]
\[
  \xi_1, \ldots, \xi_n \qquad c_1(N), \ldots, c_d(N)
\]
\[
  \Bigl( \prod_{1 \leqslant i \leqslant n} (1 + \xi_i) \Bigr)
  \sum_{0 \leqslant j \leqslant d} c_j(N) = \sum_{0 \leqslant h \leqslant
  N} c_h(E)
\]
\note{un exposant de $c_i$, illisible (\ill{}), est omis ; cet exposant et l'indice sous le dernier $\Sigma$ de la
première ligne sont peu lisibles ; en bas de cette moitié, deux lettres
isolées, lues \uncertain{$E$} et \uncertain{$N$}}

\page{66}
\note{feuillet double à l'italienne ; la moitié gauche, en partie tête-bêche,
porte des calculs de brouillon épars, dont on lit : « $e \to I^n$,
$(\varepsilon_1, \ldots, \varepsilon_n)$, $(\varepsilon_1, \ldots,
\varepsilon_n, x_{n+1}, \ldots, x_{n+p})$ » ; une flèche $I^{n+p} \to I$
marquée $x_1, \ldots, x_n$ ; « $I^p \to I^{n+p}$ » au-dessus de
« $\bigl(\prod_{i \leqslant n} (\varepsilon_i \times I)\bigr)^p \times
\ldots \subset (I \times I)^{n \times p} \times I^n$ » ; encadré, «
$\varepsilon_i \times I \subset I \times I$ » au-dessus de « $e \to I
\times I$ », avec « $(x, y)$ » et « $(x, xy)$ » ; « $f_1, \ldots, f_n$ \quad
$g_1, \ldots, g_p$ \quad $f_i g_j$ » ; une flèche $I^{n+p} \to (I \times
I)^{np} \times I^n$ marquée $(f_i, g_j)$ ; « $(x, y)$, $(x', y')$, $(x,
xy) = (x', x'y')$ », encadré « $x = x'$ », puis « $xy = x'y'$ » et une
ligne biffée ; deux segments superposés avec une flèche ; « $a) \; y = 0$
». Le reste, écrit tête-bêche, reprend les mêmes produits $(I \times I)^{np}$
et n'est pas lisible au-delà}

\begin{tikzcd}
X \arrow[r, hook] \arrow[d] \arrow[dr] & Y \arrow[d] \\
Z \arrow[r, hook] & I^m
\end{tikzcd}

comm

\begin{tikzcd}
X \arrow[r, hook] \arrow[d] & Y \arrow[d, "{\varphi = (\varphi_1, \ldots, \varphi_n)}"] \\
e \arrow[r] & I^n
\end{tikzcd}

\note{à droite, un second carré $X \subset Y$ au-dessus de $I^p
\hookrightarrow I^{p+n}$, avec « $\varphi_1, \ldots, \varphi_{p+n}$ », est
biffé de grands traits}

\[
  Y \hookrightarrow I^p, \qquad \psi = (\psi_1, \ldots, \psi_p)
\]
\[
  (\varphi_i, \psi_j) : Y \longrightarrow (I \times I)^{n \times p}
\]
\note{sous $(\varphi_i, \psi_j)$, une accolade marquée $X_{ij}$}
\struck{$(\varphi_i, \psi_j)$ se} $X_{ij} | X$ se factorise par $e \times I
\subset I \times I$

\begin{tikzcd}
X \arrow[r, hook] \arrow[d] & Y \arrow[d] \\
X' = I^p \arrow[r, hook] \arrow[d] & Y' = I^{n+p} \arrow[d] \\
e \arrow[r, hook] & S
\end{tikzcd}

\note{les deux carrés portent chacun un signe $\times$ en leur milieu}

\[
  Y \times Y \leftarrow Y \times_S Y = \underbrace{\bigl( Y \times_{Y'} Y
  \bigr)}_{\Delta_{Y/Y'}} \cup (X \times X)
\]
\note{il écrit d'abord, puis biffe, « $= Y \times_{Y'} Y = Y$ » ; l'accolade
marquée $\Delta_{Y/Y'}$ est au-dessus}
\[
  Y' \times Y' \supset Y' \times_S Y' \neq \Delta_{Y'} \cup (X' \times X')
\]
\note{le signe $\neq$ est une lecture douteuse (\uncertain{$\neq$}) ; au-dessous de
$\Delta_{Y'}$, des guillemets}
\[
  I^p \text{ --- } I^{n+p}
\]
\note{dessous, une ligne biffée, illisible}

\page{67}
\begin{tikzcd}
X \arrow[r, hook, "i"] \arrow[d, "p"'] & Y \arrow[d, "?"] \\
Z \arrow[r, hook, "?"] & I^n
\end{tikzcd}

avec a) b) c)\note{les deux points d'interrogation sont de sa main, à la
place des noms de flèches}

$f_1, \ldots, f_n : Z \to I$ \qquad $Z \to I^n$ mono

$f_\alpha \circ p = g_\alpha \circ i \,/_{1 \leqslant \alpha \leqslant p}$
\quad $Y \to I$\note{les lettres de $f_\alpha$ et $g_\alpha$ sont
surchargées ; lecture douteuse}

OK pour a)

\[
  Z \overset{\sigma}{\longrightarrow} I^q
\]
$f_{p+1}, \ldots, f_{p+q} = \sigma$

$f_{p+\beta} \circ p = \uncertain{\dot{0}}$

$g_{p+\beta} \circ i = \dot{0}$ \qquad $\bigcap_{1 \leqslant \beta
\leqslant q} g_{p+\beta}^{-1}(0) = X$ \ill{}

OK pour b)

comme

\begin{tikzcd}
X \arrow[r, hook] \arrow[d] & Y \arrow[d] \\
e \arrow[r] & I^q
\end{tikzcd}

cart.
\[
  Y \times_{I^q} Y = \Delta_Y \cup (X \times X)
\]

\begin{tikzcd}
X \arrow[r, hook, "i"] \arrow[d, "p"'] & Y \arrow[d] \\
Z \arrow[r] \arrow[d] & I^{p+q} = S \arrow[d] \\
e \arrow[r] & I^q = S'
\end{tikzcd}

\[
  \begin{array}{ccc}
    Y \times_S Y & \hookrightarrow & S \times_S S \\
    \updownarrow & & \updownarrow \delta \\
    Y \times_{S'} Y & \hookrightarrow & S \times_{S'} S \\
    \parallel & & \\
    \Delta_Y \cup X \times X & &
  \end{array}
\]
\note{schéma redessiné ; les objets de la ligne du haut sont surchargés,
on lit à droite « $S \times_S S$ » sur un premier essai biffé}

\page{69}
\note{il numérote le feuillet « 1) » en haut à gauche}

1. Soit $R$ un ensemble (ce sera bientôt $\mathbb{R}$). \struck{\ill{}}
Soit pour \struck{tout} $n \in \mathbb{N}$, $C_n \in R^n$\note{sic ; plus
bas il écrit $C_I \subset \mathfrak{P}(R^I)$}. On suppose
\begin{itemize}
  \item[(1)] $C_n$ stable par $\mathfrak{S}_n$
  \item[(2)] Soient $p, n \in \mathbb{N}$, $p \leqslant n$, considérons
  $\varphi_{p,n} : R^n \to R^p$, $(x_1, \ldots, x_n) \mapsto (x_1,
  \ldots, x_p)$, alors $\varphi_{p,n}(X) \in C_p$ si $X \in C_n$
  \item[(3)] Soient $p_1, p_2, \ldots, p_n \in \mathbb{N}$, \add{$p_{*} =
  (p_1, \ldots, p_n)$,} considérons
  \[
    R^n \xrightarrow{\ \varphi_{p_{*}}\ } R^{p_1 + \cdots + p_n}, \quad
    (x_1, \ldots, x_n) \mapsto (\underbrace{x_1, \ldots,
    x_1}_{p_1\ \text{fois}}, \ldots, \underbrace{x_n, \ldots,
    x_n}_{p_n\ \text{fois}})
  \]
  \struck{d) Il y a tel que $C_n \neq \emptyset$ ($\Longleftrightarrow$
  $C_0 \neq \emptyset$)}

  alors $\varphi_{p_{*}}(X) \in C_p$ si $X \in C_n$
  \item[(4)] Si $X \in C_m$, $Y \in C_n$, alors $X \times Y \in C_{m+n}$
\end{itemize}
\note{les numéros (1) à (4) sont cerclés ; au-dessus de $R^{p_1 + \cdots +
p_n}$, une accolade marquée $p$}

\emph{NB} Il revient au même, pour tout ens.\ fini $I$, de se donner
\struck{\ill{}} $C_I \subset \mathfrak{P}(R^I)$, de telle façon que pour
tout $\alpha : I \to J$ ($I$, $J$ ens.\ finis) d'où $\alpha^{*} : R^J \to
R^I$, on ait $\alpha^{*}(X) \in C_I$ pour $X \in C_J$, et que $X \in C_I$,
$Y \in C_J \Rightarrow X \times Y \in C_{I \sqcup J}$.

\emph{Prop} Si $X, Y \in C_n$, alors $X \cap Y \in C_n$.

\emph{Déf} Soit $X \in C_I$, $Y \in C_J$, $f : X \to Y$, d'où $\Gamma_f
\subset X \times Y \subset R^{I \sqcup J}$. On dit que $f$ est un
$C_{*}$-morphisme si $\Gamma_f \in C_{I \sqcup J}$.\note{le signe accolé à
$C$ dans « $C$-morphisme », « $C$-modèle », tout au long des pages 69 à
72, est une petite marque en indice, étoile ou croix ; transcrit $C_{*}$
partout où il figure, $C$ là où il manque} [NB.\ ne dépend que du composé
$X \to Y \hookrightarrow R^J$]\note{sous la formule, un mot biffé}

\struck{\emph{Prop} Soit $X, Y \in C_n$, alors $X \cap Y \in C_n$}

\emph{Prop} a) Si $X \in C_I$, $\mathrm{id}_X$ est \struck{\ill{}}
$C$-morphisme.

b) Le composé de deux $C$-morphismes est \uncertain{aussi un}\note{fin de
ligne serrée contre le bord}.

c) Soit $f : X \to Y$ \struck{\ill{}} une bijection. Pour que $f$
\add{$f^{-1}$ soit} \ill{} aussi un $C_{*}$-morphisme, il f.\ et s.\ que
\ill{}\note{l'addition « $f^{-1}$ soit » est en exposant au bout de la
ligne ; la construction de la phrase reste incertaine}

d) Soient $X \in C_I$, $Y \in C_J$, $Z \in C_K$, $g : X \to Y$, $h : X \to
Z$, d'où $f = (g, h) : X \to Y \times Z$ \struck{\ill{}} (où $Y \times Z
\in C_{J \sqcup K}$). Alors $f$ est un $C$-morphisme ssi $g$ et $h$ le
sont.

\marginal{e) Soit $\alpha : I \to J$, d'où $\alpha^{*} : R^J \to R^I$,
soit $X \in C_J$, \uncertain{alors} $Y = \alpha^{*}(X) \in C_I$. Alors
$\alpha^{*} | X : X \to Y$ est un $C_{*}$-morphisme.
f) Si $X \subset Y$, $X, Y \in C_I$, alors l'inj.\ $X \hookrightarrow Y$
est un $C_{*}$-morphisme.}\note{en biais dans le coin inférieur gauche ;
dans e), « $\alpha : I \to J$ » est écrit sur une première formule biffée}

\marginal{g) Si $f : X \to Y$ \ ($X \in C_n$, $Y \in C_m$) est un
$C_{*}$-morphisme, alors $f(X') \in C_m$ pour $X' \subset X$, $X' \in
C_n$.
h) Si $f$ \struck{\ill{}} comme dessus, alors $f^{-1}(Y') \in C_n$ si $Y'
\subset Y$, $Y' \in C_m$.}\note{écrit verticalement le long de la marge
gauche ; les lettres g) et h) sont peu sûres, les indices $n$, $m$ aussi}

\emph{Définition}\note{le mot est dans la marge, souligné} \struck{\ill{}}
Une structure \add{$T$ de $C_{*}$-modèle} sur un ens.\ $E$ est la donnée,
pour \struck{tout} $n \in \mathbb{N}$, d'un ens.\ $T_n$ d'applications
injectives de $E$ dans $R^n$, telles que

\page{70}
\begin{itemize}
  \item[a)] $\exists\, n \in \mathbb{N}$ tel que $T_n \neq \emptyset$
  \item[b)] Si \struck{$u \in T_n$} $(u : E \to R^n) \in T_n$, et si $v : E
  \hookrightarrow R^m$, alors $v \in T_m$ ssi 1°) $v(E) \in C_m$ \quad
  2°) $v u^{-1} : u(E) \to v(E)$ est un $C_{*}$-morphisme \quad ($u(E)
  \subset R^n$, $v(E) \subset R^m$)
\end{itemize}
\note{en haut à droite, une variante : « b) $\forall m \in \mathbb{N}$ et
$v \in T_m$, \ldots\ $v(E) \in C_m$ »}

\marginal{Cette structure peut aussi être définie par l'ensemble $\Sigma$
des $C_{*}$-morphismes de $E$ dans $R$, avec l'axiome suivant :
a) $\exists\, J \subset \Sigma$, partie finie, telle que $E \to R^J$ soit
injectif, et \uncertain{ait} image $X \in C_J$.
b) On peut choisir $J$ tel que pour $f : E \to R$, on ait $f \in \Sigma$
ssi le composé $X \simeq E \xrightarrow{f} R$ est $C_{*}$-morphisme,
[l'inverse de $E \xrightarrow{\sim} X$].}\note{en biais dans la marge
gauche, depuis le haut ; dans b), « $E \to R$ » est écrit sur une première
lettre biffée ; au-dessous, d'autres fragments en biais, « $X \subset$ »,
« $E \xrightarrow{f} R$ », non lus davantage}

\emph{Ex} Si $n \in \mathbb{N}$, \struck{$X \in R^n$} \ill{} $X \subset
C_n$\note{sic ; on attend « $X \in C_n$ »}, alors $X$ est muni d'une
$C_{*}$-structure canonique.

\emph{Déf} Soient $X, Y$ deux $C_{*}$-\struck{modèles} modèles. Une
application $f : X \to Y$ est dite un $C_{*}$-morphisme si \ldots

\emph{Prop} Les \struck{cat.} $C_{*}$-modèles, et $C_{*}$-morphismes,
forment une catégorie (pour la composition ensembliste des morphismes).
Dans cette catégorie, les produits \struck{finis} de deux objets existent.
Un morphisme de $C_{*}$-modèles est un iso ssi il est bijectif.

\emph{Prop} C.\ \uncertain{équivalentes} sur $C$
\begin{itemize}
  \item[a)] \struck{\ill{}} $\exists\, n \in \mathbb{N}$ et $X \in C_n$,
  tel que $X \neq \emptyset$
  \item[b)] $R^0 \in C_0$
  \item[c)] \struck{La} catégorie des $C_{*}$-modèles admet des objets de
  cardinal 1 (\ill{} \ill{} \ill{} objets de cardinal 1 \uncertain{sont}
  des structures de $C_{*}$-modèle) \ill{}, \ill{} \ill{} \ill{}
\end{itemize}
Alors les \struck{\ill{}} $C_{*}$-produits \struck{\ill{}} \ill{} des 2
\ill{} \ill{} objets \ill{} de la catégorie des $C_{*}$-modèles.

\emph{Dorénavant}, on suppose cette condition, sous la forme plus forte
$\forall n \in \mathbb{N}$, on a
\[
  \text{(5)} \qquad R^n = \bigcup_{X \in C_n} X
\]
\note{le 5 est cerclé et écrit sur une première formule biffée}

On définit la notion de $C_{*}$-partie d'un $C_{*}$-modèle, de façon
évidente, de façon que (i) \uncertain{compatible} « transport de structure
» et (ii) si $X \in C_n$ ($n \in \mathbb{N}$), alors pour $Y \subset X$, $Y$
est une $C_{*}$-partie du $C_{*}$-modèle $X$, ssi $Y \in C_n$. Ceci dit,
pour un morphisme de modèles, l'image directe ou inverse d'une
$C_{*}$-partie est itou.\marginal{Notion de $C_{*}$-str.\
induite}\note{séparée du texte par un trait oblique, en bas à droite ;
« induite » est lu sous réserve}

\page{71}
2. Je suppose maintenant que $R$ est un espace top.\ \add{séparé}, et je
suppose \add{$n \in \mathbb{N}$,}
\begin{itemize}
  \item[(7)] Si $X \in C_n$, $X$ est \emph{compact}.
\end{itemize}
\note{le 7 est cerclé, écrit sur un autre chiffre}

Cela implique que si $X \in C_m$, $Y \in C_p$, tout $C_{*}$-morphisme de $X$
dans $Y$ est \emph{continu}, que les $C_{*}$-modèles s.t de façon
\uncertain{naturelle} munis d'une structure d'espace compact, d'où un
foncteur
\[
  (C_{*}\text{-modèles}) \longrightarrow \text{espaces compacts.}
\]

Je suppose de plus
\begin{itemize}
  \item[(8)] $\forall\, n \in \mathbb{N}$, $X \in C_n$, et $x \in X$,
  alors \struck{$V \in C_n$ qui soit un voisinage de $x$ dans $X$, $\exists
  Y$} les $Y \in C_n$ qui \add{dans $X$} sont voisinage de $x$ forment un
  syst.\ fondamental de voisinages de $x$ dans $X$.
\end{itemize}
\note{le 8 est cerclé ; les passages biffés sont récrits au-dessus de la
ligne}
[Il suffit que \struck{\ill{}} les intérieurs des $X \in C_n$ forment une
base de la topologie de $R^n$]\marginal{\ill{} \ill{} \ill{}}\note{trait
vertical dans la marge gauche, contre le crochet, avec trois mots en biais,
illisibles}

Alors, les \add{$C_{*}$-}morphismes d'un $C_{*}$-modèle $X$ dans un autre
$Y$ sont les sections d'un faisceau --- le faisceau des germes de
$C_{*}$-morphismes de $X$ dans $Y$. Il en est ainsi pour les
$C_{*}$-morphismes de $X$ dans $R$ (\ill{} dans un $R^I$). La connaissance
du faisceau des $C_{*}$-morphismes d'un modèle $X$ \add{(comme $H$-faisceau
\ill{} \ill{} \ill{} \ill{} \ill{} $X$ dans $R$)} dans $R$, ne donne
\uncertain{pas} la structure de modèle.\note{addition interlinéaire
serrée, lue en partie}

$C_{*}$-espace : espace loc.$^t$ compact, muni d'un faisceau \ill{}
\uncertain{faisceau} des \uncertain{fonctions} continues \ill{} \ill{} \ill{}
\ill{} de $R$, tel que \ill{} pt ait un voisinage \uncertain{compact}
\struck{\ill{}} qui, pour la structure induite, soit un $C_{*}$-modèle.

\page{72}
\note{il numérote le feuillet « 2) » en haut à gauche}

Considérons la condition
\begin{itemize}
  \item[(6)] \struck{$\forall$} $\forall n \in \mathbb{N}$, $C_n$ stable
  par réunions finies \quad ($\Rightarrow \emptyset \in C_n$)
\end{itemize}
\note{le 6 est cerclé}

On en conclut

\emph{Prop} Soit $X$ un $C_{*}$-modèle, $(X_i)_{i \in I}$ une famille finie
de $C_{*}$-parties de $X$, avec $X = \bigcup X_i$, $f : X \to Y$ une
application de $X$ dans un $C_{*}$-modèle $Y$. Alors $f$ est un
$C_{*}$-morphisme ssi $\forall i \in I$, $f | X_i$ l'est.

\emph{Corollaire} Dans la cat.\ des $C_{*}$-modèles, $X$ est la somme
amalgamée des $X_i$, sur les $X_i \cap X_j$.

\emph{Question} Soit un diagramme

\begin{tikzcd}[column sep=small]
 & Z \arrow[dl, hook'] \arrow[dr, hook] & \\
X & & Y
\end{tikzcd}

dans la catégorie des $C_{*}$-modèles, y a-t-il une somme amalgamée ? OK.
Si on \struck{\ill{} \ill{}} \add{\ill{} la condition}
suivante : $\forall n \in \mathbb{N}$, $X, Y \in C_n$, $Y \subset X$, et
$\forall$ $C_{*}$-morphisme $f_Y : Y \to R$, $\exists$ $C_{*}$-morphisme
$f_X : X \to R$ qui le prolonge.\marginal{? \ \ill{} \ill{}}\note{dans la
marge gauche, un point d'interrogation et deux mots, dont le premier se
termine peut-être en « -tion », à côté d'un trait vertical qui borde la
condition}

\end{document}
